\section{Диаграммы последовательности}
\headingtextsep

Диаграммы последовательности используются для описания поведения системы во времени и показывают порядок обмена сообщениями между участниками сценария. Для системы «Контур» выбраны три ключевых сценария: основной пользовательский путь, альтернативный сценарий отказа в доступе и сложный интеграционный сценарий с внешними системами. Соответствующие диаграммы приведены в приложении А на рисунках \ref{fig:appendix-sequence-happy-path}, \ref{fig:appendix-sequence-access-denied} и \ref{fig:appendix-sequence-integration-sync}.

\textheadingsep
\subsection{Основной сценарий: создание задачи и фиксация связанного контекста}
\headingtextsep

В качестве основного сценария выбран процесс создания задачи в проекте. Этот путь является типовым для координатора или исполнителя и напрямую связан с центральной функцией системы --- управлением работами. В сценарии участвуют веб-клиент, API-сервис, модуль проверки доступа, сервис управления задачами, модуль связывания сущностей, PostgreSQL и асинхронный обработчик аудита.

Сценарий начинается с заполнения формы задачи в пользовательском интерфейсе и отправки синхронного HTTP-запроса в API. После проверки сессии и прав доступа сервис задач создает запись в базе данных. Если при создании задачи выбраны связанные документы, заметки или файловые ресурсы, система переходит к условной ветви и в цикле добавляет соответствующие связи. После успешного завершения синхронной части формируется асинхронное сообщение в очередь событий, по которому отдельный обработчик сохраняет запись аудита. Такой подход позволяет сразу вернуть пользователю результат создания задачи и не задерживать ответ операциями фонового журналирования.

\textheadingsep
\subsection{Альтернативный сценарий: отказ в изменении статуса задачи агентом}
\headingtextsep

Альтернативный сценарий моделирует ситуацию, в которой агент пытается изменить статус задачи без необходимого разрешения на такую операцию. В этом сценарии взаимодействуют агент, MCP-мост, API-сервис, модуль аутентификации и проверки набора разрешений, сервис задач и контур аудита. Диаграмма показывает точку возникновения ошибки, а также то, что бизнес-операция прерывается до момента изменения данных.

Запрос начинается с вызова инструмента агентом, после чего MCP-мост преобразует его в синхронный HTTP-запрос к API. Ошибка возникает на этапе проверки действительности ключа и набора разрешений. Если ключ не содержит нужного разрешения, запрос не передается в сервис изменения задачи, а вызывающая сторона получает ответ с кодом 403 Forbidden и сообщением об отказе в доступе. Для случая недействительного или истекшего ключа в диаграмме также может быть показана альтернативная ветвь с кодом 401 Unauthorized. Отказ фиксируется отдельным асинхронным событием аудита, что позволяет сохранить прослеживаемость неуспешных попыток изменения состояния задачи.

\textheadingsep
\subsection{Сложный сценарий: синхронизация артефактов GitHub и файлового хранилища}
\headingtextsep

Сложный сценарий описывает обработку задачи, связанной с внешними инженерными артефактами. Пользователь инициирует синхронизацию карточки задачи, после чего API-сервис создает асинхронное задание и передает его в очередь сообщений. Далее отдельный обработчик координирует взаимодействие как минимум с двумя внешними интеграциями: GitHub API и S3-совместимым файловым хранилищем.

Внутри сценария показано параллельное выполнение двух ветвей. В первой ветви обработчик получает по HTTP сведения о связанных issue, pull request, commit и branch из GitHub. Во второй ветви обработчик записывает в объектное хранилище сформированный снимок связанных файловых материалов и метаданные вложений. Для обеих ветвей предусматриваются таймауты и цикл повторных попыток. Если внешний вызов не завершается успешно за отведенное время, обработчик выполняет ограниченное число retry-повторов, после чего переводит задачу синхронизации в состояние частичного или полного отказа. После успешного завершения параллельных ветвей метаданные сохраняются в PostgreSQL, а в очередь публикуется событие для дальнейшей индексации и отражения результата в пользовательском интерфейсе.
